\chapter{Clase Taller 5}
Un volumen de fluido. Un elemento de volumen $dV$.

Sobre el elemento de volumen $dV$ hay dos tipos de fuerzas:

\begin{enumerate}
    \item Fuerzas de superficie: $\vec{F}_p$ actúa sobre la superficie.
    \item Fuerzas de volumen: $\vec{F}_e$ actúa sobre todo el volumen.
\end{enumerate}

Si el sistema está en equilibrio:
$$
d\vec{F}_e + d\vec{F}_p = 0
$$
con
$$
d\vec{F}_p = -p\,d\vec{S}
$$

De la integral:
$$
-\oint p\,d\vec{S} + \int \vec{f}_e\,dV = 0
$$

Por Gauss:
$$
\oint p\,d\vec{S} = \int \nabla p\,dV
$$

Entonces:
$$
\boxed{\nabla p = \vec{f}_e}
$$
Ec.\ de equilibrio (en general).

Para el caso gravitacional:
$$
\nabla p = \rho\,\vec{g}
$$
(en coordenadas).

De sig.\ que $\nabla \times \vec{f} = 0$ $\Rightarrow$ fuerzas conservativas.

\vspace{1em}

\noindent\textbf{Principio de Pascal:} la presión aplicada sobre un fluido en un contenedor se transmite por igual a cada porción del fluido y a las paredes del recipiente.

\vspace{1em}

\noindent\textbf{Ecuación de estado:} relación entre las variables termodinámicas.
$$
F(p,\rho,T) = 0
$$
Para un gas ideal:
$$
pV = nRT \quad \Longrightarrow \quad p = \rho R T
$$
(con $R$ la constante específica del gas).

La ec.\ de equilibrio también se puede escribir como:
$$
\nabla\Phi + \frac{\nabla p}{\rho} = 0
\quad\Longrightarrow\quad
\Phi^* = \Phi + w = \text{const}
$$

Definiendo el potencial de presión y el potencial efectivo:
$$
\Phi^* = \Phi + w = \text{const}
$$
con
$$
w = \int \frac{dp}{\rho}
$$

Para un fluido politrópico:
$$
p = C\,\rho^{\gamma}
$$

Módulo de compresibilidad:
$$
K = \rho\,\frac{dp}{d\rho}
$$

\vspace{1em}

\noindent\rule{\textwidth}{0.4pt}

\vspace{1em}

\noindent{\Large\bfseries Medios continuos autogravitantes}

\vspace{1em}

En general, el campo gravitacional de una distribución de masa:
$$
\Phi(\vec{r}) = -G\int \frac{\rho(\vec{r}\,')}{|\vec{r}-\vec{r}\,'|}\,d^3\vec{r}\,'
$$

Por Gauss:
$$
\oint \vec{g}\cdot d\vec{S} = -4\pi G M
$$

En coordenadas:
$$
\vec{g} = -\nabla\Phi
$$

Por tanto:
$$
\boxed{\nabla^2\Phi = 4\pi G\rho}
\qquad \text{(Poisson)}
$$

De la ec.\ de equilibrio se sigue que $\nabla p = -\rho\,\nabla\Phi$.

\vspace{1em}

\noindent\textbf{Modelo de la estructura estelar.}

La ec.\ de equilibrio:
$$
\nabla p = -\rho\,\nabla\Phi
$$

Ec.\ barotrópica. Una ecuación de estado conveniente para la politropa:
$$
p = K\,\rho^{1+\frac{1}{n}}
$$

Combinando con Poisson:
$$
\nabla^2\Phi = 4\pi G\rho
$$

Haciendo $\rho = \rho_c\,\theta^n$, $\xi = \left(\dfrac{4\pi G\rho_c}{\Phi_T - \Phi_c}\right)^{1/2}r$, se obtiene la ec.\ de Lane--Emden:
$$
\boxed{\frac{1}{\xi^2}\frac{d}{d\xi}\!\left(\xi^2\frac{d\theta}{d\xi}\right) = -\theta^n}
$$

\vspace{1em}

\noindent\textbf{Condiciones de frontera:}

\begin{itemize}
    \item $\theta(0) = 1$: por definición.
    \item $\theta'(0) = 0$: fuerza nula en $r = 0$.
\end{itemize}

\vspace{1em}

\noindent\textbf{Soluciones:}

\begin{itemize}
    \item $n = 0$: solución directa
    $$
    \theta(\xi) = 1 - \frac{\xi^2}{6}
    $$
    Superficie del astro: $\theta = 0$, es decir $\xi_1 = \sqrt{6}$.

    \item $n = 1$:
    $$
    \theta(\xi) = \frac{\sin\xi}{\xi}
    $$
    Superficie del astro: $\theta = 0$, es decir $\xi_1 = \pi$.

    \item $n = 5$:
    $$
    \theta(\xi) = \left(1 + \frac{\xi^2}{3}\right)^{-1/2}
    $$
    Superficie: $\theta \to 0$ cuando $\xi \to \infty$.
\end{itemize}